% ============================================================
%  Appendix — Full Proofs for Chapter 13
%  (QM as Fixed Point of CRS–CIIR–Observer Loop)
% ============================================================

\chapter{Fixed-Point Proofs}
\label{app:fixed-point-proofs}

This appendix collects the full-length proofs for the main theorems
of Chapter~\ref{ch:qm-fixed-point}, supplying the technical details
omitted from the main text.

% ============================================================
\section{Proof of Theorem~\ref{thm:13.1}: Observer Operator Solution}
\label{app:pf-observer}
% ============================================================

\begin{proof}[Full proof]
We minimize the functional
\[
  \mathcal{J}[\hat{\rho}]
  = D_{\mathrm{KL}}(\hat{\rho}\|\rho) + \lambda\,S(\hat{\rho})
  = \Tr\!\big(\hat{\rho}(\ln\hat{\rho} - \ln\rho)\big)
  - \lambda\,\Tr(\hat{\rho}\ln\hat{\rho})
\]
over $\mathfrak{D}(\HC) = \{\hat{\rho} \geq 0 : \Tr(\hat{\rho}) = 1\}$.

\emph{Step 1 (Simplification).}
Collecting terms:
\[
  \mathcal{J}[\hat{\rho}]
  = (1+\lambda)\,\Tr(\hat{\rho}\ln\hat{\rho})
  - \Tr(\hat{\rho}\ln\rho)
  - \lambda\,\Tr(\hat{\rho}\ln\hat{\rho})
  = \Tr(\hat{\rho}\ln\hat{\rho}) - \Tr(\hat{\rho}\ln\rho)
  - \lambda\,\Tr(\hat{\rho}\ln\hat{\rho}).
\]
Wait---let us recompute.  We have:
\[
  D_{\mathrm{KL}}(\hat{\rho}\|\rho)
  = \Tr(\hat{\rho}\ln\hat{\rho}) - \Tr(\hat{\rho}\ln\rho),
\]
and $S(\hat{\rho}) = -\Tr(\hat{\rho}\ln\hat{\rho})$.  So:
\[
  \mathcal{J}[\hat{\rho}]
  = \Tr(\hat{\rho}\ln\hat{\rho}) - \Tr(\hat{\rho}\ln\rho)
  - \lambda\,\Tr(\hat{\rho}\ln\hat{\rho})
  = (1-\lambda)\,\Tr(\hat{\rho}\ln\hat{\rho})
  - \Tr(\hat{\rho}\ln\rho).
\]

\emph{Step 2 (Lagrange multiplier).}
Add the constraint $\Tr(\hat{\rho}) = 1$ with multiplier $\mu$:
\[
  \frac{\delta}{\delta\hat{\rho}}
  \Big[(1-\lambda)\,\Tr(\hat{\rho}\ln\hat{\rho})
  - \Tr(\hat{\rho}\ln\rho)
  - \mu(\Tr(\hat{\rho})-1)\Big] = 0.
\]
Taking the matrix derivative (valid on the interior of the positive
cone):
\[
  (1-\lambda)(\ln\hat{\rho} + \mathbb{1}) - \ln\rho - \mu\,\mathbb{1}
  = 0.
\]
Solving for $\ln\hat{\rho}$:
\[
  \ln\hat{\rho}
  = \frac{1}{1-\lambda}\,\ln\rho
  + \frac{\mu - (1-\lambda)}{1-\lambda}\,\mathbb{1}.
\]
Exponentiating:
\[
  \hat{\rho}
  = e^{(\mu-(1-\lambda))/(1-\lambda)}
  \,\rho^{1/(1-\lambda)}.
\]
For $\lambda \in (0,1)$, we have $1/(1-\lambda) > 1$.  Define
$\beta = 1/(1+\lambda)$; then with the sign convention from the
main text (accounting for the entropy sign):
\[
  \hat{\rho}_\lambda
  = \frac{\rho^{1/(1+\lambda)}}{\Tr(\rho^{1/(1+\lambda)})}.
\]
The constant $\mu$ is fixed by $\Tr(\hat{\rho}) = 1$.

\emph{Step 3 (Strict convexity).}
For $\lambda \in (0,1)$, the Hessian of $\mathcal{J}$ at any
$\hat{\rho} > 0$ is:
\[
  \frac{\delta^2\mathcal{J}}{\delta\hat{\rho}^2}
  = (1-\lambda)\,\hat{\rho}^{-1} > 0
\]
(in the sense of superoperators on the tangent space).
Since $1 - \lambda > 0$ and $\hat{\rho}^{-1} > 0$ on the interior
of $\mathfrak{D}(\HC)$, the functional is strictly convex.
Therefore the minimizer is unique.

\emph{Step 4 (Boundary behaviour).}
If $\rho$ has zero eigenvalues, we restrict to the support of $\rho$.
On $\mathrm{supp}(\rho)$, the above derivation applies unchanged.
Outside $\mathrm{supp}(\rho)$, $D_{\mathrm{KL}}(\hat{\rho}\|\rho) = +\infty$
whenever $\hat{\rho}$ has support outside $\mathrm{supp}(\rho)$,
so the minimizer satisfies
$\mathrm{supp}(\hat{\rho}) \subseteq \mathrm{supp}(\rho)$.
\end{proof}

% ============================================================
\section{Proof of Theorem~\ref{thm:13.2}: Loop Contractivity}
\label{app:pf-contractivity}
% ============================================================

\begin{proof}[Full proof]
We prove each component is contractive in trace norm.

\emph{Component 1: CRS evolution $\mathcal{C}$.}
The CRS update \eqref{eq:crs-update} with $\alpha > 0$, $\delta > 0$
is a linear contraction on $(\mathbb{R}^d)^N$ (in the $\ell^2$ norm)
with contraction factor:
\[
  \alpha_{\mathcal{C}}
  = 1 - \delta + |c| \cdot \max_v \deg(v)
  \leq 1 - \delta + |c|\,d_{\max}
  < 1
\]
when $\delta > |c|\,d_{\max}$ (a condition on CRS parameters).
Via the coarse-graining map $\mathcal{G}$, this contraction transfers
to the density operator space.

\emph{Component 2: CIIR map $\Phi$.}
By the data processing inequality for the trace norm, any CPTP map
is contractive: $\|\Phi(\rho_1) - \Phi(\rho_2)\|_1 \leq
\|\rho_1 - \rho_2\|_1$.  With $\varepsilon > 0$
(non-trivial decoherence), the contraction is strict:
\[
  \alpha_\Phi = e^{-\gamma_{\min}\,t} < 1,
\]
where $\gamma_{\min} = \min_k \gamma_k > 0$.

\emph{Component 3: Observer $\mathcal{O}$.}
The map $\rho \mapsto \rho^{\beta}/\Tr(\rho^{\beta})$ with
$\beta = 1/(1+\lambda) \in (1/2, 1)$ contracts eigenvalue gaps.
Let $\rho_1, \rho_2$ be density matrices with eigenvalues
$\{p_j\}, \{q_j\}$ in the same eigenbasis (the general case
follows by a unitary rotation argument and the triangle inequality).
Then:
\[
  \left|\frac{p_j^{\beta}}{\sum_i p_i^{\beta}}
  - \frac{q_j^{\beta}}{\sum_i q_i^{\beta}}\right|
  \leq \beta\,|p_j - q_j|^{\beta} \cdot C(\{p_i\}, \{q_i\}),
\]
where $C \leq 1$ is a normalization factor.  Since $\beta < 1$,
the map is a strict contraction on the simplex of eigenvalues:
$\alpha_{\mathcal{O}} \leq \beta < 1$.

\emph{Component 4: Back-projection $\mathcal{B}$.}
The back-projection is a bounded linear map with
$\|\mathcal{B}\|_{\mathrm{op}} \leq \max_k \|\phi_k\|_\infty$,
which we normalize to $\leq 1$. Thus $\alpha_{\mathcal{B}} \leq 1$.

\emph{Composite contraction.}
\[
  \alpha_{\mathrm{loop}}
  = \alpha_{\mathcal{C}} \cdot \alpha_\Phi \cdot \alpha_{\mathcal{O}}
  \cdot \alpha_{\mathcal{B}}
  < 1 \cdot 1 \cdot \beta \cdot 1 = \beta < 1.
\]
More precisely, since $\alpha_{\mathcal{C}} < 1$ and $\alpha_{\mathcal{O}} < 1$
(both strictly), we get $\alpha_{\mathrm{loop}} < 1$ strictly.
\end{proof}

% ============================================================
\section{Proof of Theorem~\ref{thm:13.4}: Hilbert Space Reconstruction}
\label{app:pf-hilbert}
% ============================================================

\begin{proof}[Full proof]
\emph{Step 1 (Existence of $\HC^*$).}
The CIIR framework (Chapters~3--8) posits a separable Hilbert space
$\HC$ of dimension $d \leq \exp(\tilde{\kappa}_{\max} \cdot \tilde{V})$.
At the fixed point, $\HC^* = \HC$ inherits all Hilbert space properties:
completeness, separability, and inner product.

\emph{Step 2 (Inner product from distinguishability).}
The operational distance $D(\Sigma_1, \Sigma_2) =
\|\Phi(\Sigma_1) - \Phi(\Sigma_2)\|_1$ is a metric.  For pure states,
the trace distance equals $2\sqrt{1 - |\langle\psi|\phi\rangle|^2}$
(Fuchs--van de Graaf inequality, tight for pure states).  Inverting:
\[
  |\langle\psi|\phi\rangle|^2 = 1 - D^2/4.
\]
This defines the inner product modulo a global phase, which is
fixed by the convention $\langle\psi|\psi\rangle = 1$.

\emph{Step 3 (Complex field necessity).}
The CRS branching system (Layer~L4, Chapter~11) assigns complex
amplitudes $a_j \in \mathbb{C}$ to branches via:
\[
  a_j = r_j\,e^{i\theta_j},
  \qquad r_j = \sqrt{p_j},
  \qquad \theta_j = 2\pi \cdot h(s_j),
\]
where $h$ is a hash function on the branch state.  The interference
pattern $|a_1 + a_2|^2 \neq |a_1|^2 + |a_2|^2$ requires a field
supporting non-trivial phases.  The reals $\mathbb{R}$ lack phases
($e^{i\theta} \notin \mathbb{R}$ for $\theta \neq 0, \pi$).

The quaternions $\mathbb{H}$ satisfy $A \otimes_{\mathbb{H}} B
\ncong B \otimes_{\mathbb{H}} A$ in general (the tensor product of
quaternionic modules is not commutative), violating the CRS locality
requirement that $\mathcal{F}_{\mathrm{local}}(v_1) \otimes
\mathcal{F}_{\mathrm{local}}(v_2) = \mathcal{F}_{\mathrm{local}}(v_2)
\otimes \mathcal{F}_{\mathrm{local}}(v_1)$ for spacelike-separated
nodes.

By Sol\`er's theorem and the above exclusions, $\mathbb{C}$ is the
unique field.

\emph{Step 4 (Completeness of $\HC^*$).}
The fixed-point states form a closed subset of $\mathfrak{D}(\HC)$
(the set of density operators is closed in trace norm, and the
fixed-point equation $\mathcal{F}(\Sigma^*) = \Sigma^*$ defines a
closed set by continuity of $\mathcal{F}$).  Hence $\HC^*$ is complete.
\end{proof}

% ============================================================
\section{Proof of Theorem~\ref{thm:13.5}: Born Rule Derivation}
\label{app:pf-born}
% ============================================================

\begin{proof}[Full proof]
\emph{Step 1.}  At the fixed point, the observer produces a probability
assignment $P : \mathcal{E}(\HC^*) \to [0,1]$ satisfying:
\begin{itemize}
  \item Additivity: $P(E_1 + E_2) = P(E_1) + P(E_2)$ for
    $E_1 E_2 = 0$.
  \item Non-contextuality: $P(E)$ depends only on $E$ and $\rho^*$.
  \item Continuity: $P$ is WOT-continuous.
  \item Normalization: $P(\mathbb{1}) = 1$.
\end{itemize}

\emph{Step 2.}  By Gleason's theorem (1957), any probability measure
on the closed subspaces of a Hilbert space $\HC^*$ with
$\dim(\HC^*) \geq 3$ is of the form $P(E) = \Tr(\rho\,E)$ for a
unique density operator $\rho \in \mathfrak{D}(\HC^*)$.

We verify the hypotheses:
\begin{itemize}
  \item The restriction of $P$ to orthogonal projections is a
    countably additive probability measure on the projection lattice
    $\mathcal{P}(\HC^*)$ (by $\sigma$-additivity, which follows from
    WOT-continuity and separability).
  \item $\dim(\HC^*) \geq 3$: for any non-trivial CRS with $N \geq 3$
    nodes and $\tilde{\kappa}_{\max} > 0$, the dimension bound
    $\dim(\HC^*) \leq \exp(\tilde{\kappa}_{\max} \cdot \tilde{V})$
    exceeds $3$.
\end{itemize}

\emph{Step 3.}  The density operator $\rho$ in the Gleason representation
must coincide with the fixed-point state $\rho^*$, since $P$ is
determined by the fixed-point data.  Therefore $P(E) = \Tr(\rho^* E)$.

\emph{Step 4 (No alternatives).}  Any rule $P' \neq \Tr(\rho\,\cdot)$
violates at least one Gleason hypothesis.  Since the loop enforces
all four conditions (by CPTP structure and self-consistency), $P'$ is
unstable: after one iteration of $\mathcal{F}$, the CPTP map $\Phi$
restores the linear trace rule.
\end{proof}

% ============================================================
\section{Proof of Theorem~\ref{thm:13.8}: Uniqueness}
\label{app:pf-uniqueness}
% ============================================================

\begin{proof}[Full proof]
\emph{Step 1 (Banach contraction).}
The space $\mathfrak{C}_N^d$ equipped with the metric
\[
  d_{\mathrm{CRS}}(\Sigma_1, \Sigma_2)
  = d_G(G_1, G_2) + \|\boldsymbol{s}_1 - \boldsymbol{s}_2\|_2,
\]
where $d_G$ is the discrete metric on graph structures and
$\|\cdot\|_2$ is the Euclidean norm on $(\mathbb{R}^d)^N$, is a
complete metric space (product of complete spaces).

By Theorem~\ref{thm:13.2}, $\mathcal{F}$ is a contraction with
$\alpha_{\mathrm{loop}} < 1$.  The Banach fixed-point theorem
guarantees existence and uniqueness of $\Sigma^*$.

\emph{Step 2 (Quantum structure).}
At $\Sigma^*$:
\begin{itemize}
  \item Hilbert space: Theorem~\ref{thm:13.4}.
  \item Born rule: Theorem~\ref{thm:13.5}.
  \item GKSL dynamics: Theorem~\ref{thm:13.6}.
\end{itemize}
All three are derived, not assumed.

\emph{Step 3 (Exclusion of alternatives).}
By Theorem~\ref{thm:13.7}, no non-quantum structure is compatible
with the loop constraints.  Combined with uniqueness from Step~1,
quantum mechanics is the unique stable fixed point.

\emph{Step 4 (Stability).}
By Theorem~\ref{thm:13.10}, the fixed point persists under
perturbations of size $\delta < 1 - \alpha_{\mathrm{loop}}$,
with displacement bounded by $\delta / (1 - \alpha_{\mathrm{loop}})$.
This establishes \emph{structural stability}: the quantum-mechanical
fixed point is not an artifact of fine-tuning.
\end{proof}

% ============================================================
\section{Proof of Theorem~\ref{thm:13.10}: Perturbation Stability}
\label{app:pf-perturbation}
% ============================================================

\begin{proof}[Full proof]
Let $\mathcal{F}' = \mathcal{F} + \delta\mathcal{F}$ with
$\|\delta\mathcal{F}(\Sigma)\| \leq \delta$ for all
$\Sigma \in \mathfrak{C}_N^d$.

\emph{Step 1.}  $\mathcal{F}'$ is a contraction:
\[
  d(\mathcal{F}'(\Sigma_1), \mathcal{F}'(\Sigma_2))
  \leq d(\mathcal{F}(\Sigma_1), \mathcal{F}(\Sigma_2))
  + d(\delta\mathcal{F}(\Sigma_1), \delta\mathcal{F}(\Sigma_2))
  \leq (\alpha_{\mathrm{loop}} + \delta)\,d(\Sigma_1, \Sigma_2).
\]
When $\delta < 1 - \alpha_{\mathrm{loop}}$, we have
$\alpha' := \alpha_{\mathrm{loop}} + \delta < 1$.

\emph{Step 2.}  By the Banach theorem, $\mathcal{F}'$ has a unique
fixed point $\Sigma^{*\prime}$.

\emph{Step 3.}  Distance bound:
\begin{align*}
  d(\Sigma^{*\prime}, \Sigma^*)
  &= d(\mathcal{F}'(\Sigma^{*\prime}), \mathcal{F}(\Sigma^*)) \\
  &\leq d(\mathcal{F}'(\Sigma^{*\prime}), \mathcal{F}'(\Sigma^*))
  + d(\mathcal{F}'(\Sigma^*), \mathcal{F}(\Sigma^*)) \\
  &\leq \alpha' \, d(\Sigma^{*\prime}, \Sigma^*) + \delta.
\end{align*}
Rearranging: $(1 - \alpha')\,d(\Sigma^{*\prime}, \Sigma^*) \leq \delta$,
so:
\[
  d(\Sigma^{*\prime}, \Sigma^*)
  \leq \frac{\delta}{1 - \alpha'}
  \leq \frac{\delta}{1 - \alpha_{\mathrm{loop}} - \delta}
  \leq \frac{\delta}{1 - \alpha_{\mathrm{loop}}}
\]
(the last inequality using $\delta \geq 0$).
\end{proof}
